\section{Inference Under Uncertainty}

So far, we have considered prediction when $z$ is directly given. We now consider prediction from a history $x_{1:t}$. The representational question remains the same: what geometry allows the task's conditional distribution to be decoded linearly? Inference changes the input to this conditional, but not the factorization principle.

Throughout, $L_c$ denotes a log-conditional centered over the variable being predicted, and $\tilde p$ denotes an unnormalized probability. We follow the convention $L_c=ED$: rows of $E$ are embedding vectors $E(s)^\top$, and columns of $D$ are decoder vectors $D(y)$.

For a prediction target $y$, define two histories as equivalent when they induce the same target distribution:
\[
x_{1:t}\sim x'_{1:t'}
\quad\Longleftrightarrow\quad
p(y\mid x_{1:t})=p(y\mid x'_{1:t'})
\qquad \forall y.
\]

Let $s_t=\Pi(x_{1:t})$ denote this equivalence class, or \emph{state}. Then
\[
p(y\mid s_t)=p(y\mid x_{1:t}),
\]
so the task kernel factors through the state:
\[
L_c(x_{1:t},y)=L_c(s_t,y).
\]

The state is target-dependent. When predicting $z$, it identifies the posterior over $z$. When predicting the next observation, it can be smaller: different posteriors may yield the same next-token distribution.

\subsection{The same factorization over histories}

We can define the task kernel directly over histories, without first identifying or enumerating its distinct states. For a sample of $N$ histories, construct
\[
L_c[i,y]
=
L_c(x^{(i)}_{1:t_i},y)
=
L_c(s_i,y).
\]

Histories inducing the same state give identical rows. We retain these repetitions: there is one row per sampled prediction instance, not one row per distinct state. Consequently, sampling histories from the task automatically represents how frequently each state is encountered.

We now apply exactly the same factorization as in the static setting:
\[
L_c=U\Sigma V^\top=ED.
\]

Repeated rows do not change the earlier gradient-flow argument. For directly optimized matrix factors, the conserved quantity remains
\[
E^\top E-DD^\top.
\]
An initialization satisfying $E^\top E=DD^\top$ therefore remains balanced; sufficiently small initial factors instead make the conserved discrepancy small. At an exact balanced solution with no redundant coordinates, the singular-value split is
\[
\boxed{
E=U\Sigma^{1/2},
\qquad
D=\Sigma^{1/2}V^\top,
}
\]
up to a common orthogonal change of embedding coordinates.

In particular, histories with the same target distribution receive the same embedding in this factorization. The state is the quotient already present in the kernel; an analytic expression for $\Pi$ is not required to derive its normative representation.

As in the static case, this balance result concerns directly optimized matrix factors. Whether the representations produced by a shared neural encoder approach the same geometry remains an empirical question.

\subsection{Computing the geometry without enumerating histories}

Although there may be infinitely many histories, the decoder geometry can be obtained from a matrix whose size depends only on the number of outputs. For an exact balanced factorization,
\[
\begin{aligned}
L_c^\top L_c
&=
D^\top E^\top E D\\
&=
D^\top(DD^\top)D\\
&=
(D^\top D)^2.
\end{aligned}
\]

Therefore,
\[
\boxed{
D^\top D
=
\left(L_c^\top L_c\right)^{1/2}
=
\left(
\sum_{i=1}^N
L_c(s_i,\cdot)^\top L_c(s_i,\cdot)
\right)^{1/2}.
}
\]

Here the square root is the positive-semidefinite matrix square root. A state appearing repeatedly contributes repeatedly to the sum. Thus its weighting follows directly from the rows of the sampled task kernel, rather than from a separate weighting assumption.

For independently sampled histories with finite logit second moments,
\[
\frac{1}{N}L_c^\top L_c
\longrightarrow
\mathbb E_{x_{1:t}}
\left[
L_c(s_t,\cdot)^\top L_c(s_t,\cdot)
\right].
\]

Equivalently, the decoder Gram, normalized for sample size, satisfies
\[
\boxed{
\frac{D^\top D}{\sqrt N}
\longrightarrow
\left(
\mathbb E_{x_{1:t}}
\left[
L_c(s_t,\cdot)^\top L_c(s_t,\cdot)
\right]
\right)^{1/2}.
}
\]

The expectation is therefore the population counterpart of ordinary SVD over sampled histories. The task distribution, including its distribution over context lengths, enters through which rows are sampled. The sample-size normalization changes only the overall scale, not the relative geometry. Centering remains over output alternatives, not across histories.

For any exactly represented state, its minimum-norm embedding is recovered from its target logits:
\[
\boxed{
E(s)
=
(D^\top)^\dagger L_c(s,\cdot)^\top.
}
\]

With fewer dimensions than the kernel's rank, retaining the leading singular modes gives the best squared-logit approximation. This is a statement about approximation of the centered log-conditional kernel, not an identification of truncated SVD with rank-constrained cross-entropy optimization.

Thus the practical procedure is the same for any prediction target: sample histories, compute the target conditionals, center their logarithms, and factorize the resulting matrix. Alternatively, accumulate its output-side Gram without storing all its rows.

\subsection{Latent inference and next-token prediction}

The two inference tasks differ in the conditional kernel being factorized, not in the representational analysis.

\paragraph{Inferring the latent variable.}

When predicting $z$, Bayes' rule supplies the target:
\[
\log\tilde p(z\mid x_{1:t})
=
\log p_0(z)+\log p(x_{1:t}\mid z).
\]
Centering over $z$ gives $L_c(s_t,z)$.

For conditionally independent observations,
\[
\log\tilde p(z\mid x_{1:t})
=
\log p_0(z)+\sum_{i=1}^t\log p(x_i\mid z).
\]

This case provides a simple illustration of the history-to-state computation. Let $E_Z$ be a basis for the normalized log-likelihood functions centered over $z$, including the centered log-prior when necessary. Denote their coordinates by $m(x)$ and $s_0$, respectively. We use $m(x)$ for the evidence message to distinguish it from the decoder $D$ of the history kernel.

These factors can be obtained from the static likelihood factorization by restoring its normalizer, centering over $z$, and retaining a basis. In particular,
\[
\log p(x\mid z)
=
\log\tilde p(x\mid z)
-
\log\sum_{x'}\tilde p(x'\mid z).
\]
The normalizer is a single function of $z$, so restoring it requires at most one additional likelihood direction. Its per-observation coefficient is constant, and its accumulated coordinate is proportional to the observation count $t$.

Centering Bayes' rule now gives
\[
L_c(s_t,\cdot)^\top
=
E_Z\left[s_0+\sum_{i=1}^t m(x_i)\right].
\]

Since $E_Z$ is a basis, we can identify the state with these unique coordinates:
\[
\boxed{
s_t=s_0+\sum_{i=1}^t m(x_i),
\qquad
L_c(s_t,\cdot)^\top=E_Zs_t.
}
\]

Importantly, additivity also holds for the balanced representation itself. Using the history-kernel decoder $D$ for the latent target,
\[
\begin{aligned}
E(s_t)
&=
(D^\top)^\dagger E_Zs_t\\
&=
E(s_0)
+
\sum_{i=1}^t
(D^\top)^\dagger E_Zm(x_i).
\end{aligned}
\]

Thus the likelihood messages and the balanced neural representation are related by a fixed linear map. The predicted representation is affine in the observation histogram and invariant to permutations of the same observations.

This prediction concerns the minimal representation, or the readout-relevant component of a larger representation. Additional components invisible to the centered readout need not be additive or permutation-invariant.

\paragraph{Predicting the next observation.}

For next-token prediction, consider
\[
x_{1:t}\leftarrow z\rightarrow x_{t+1},
\qquad
P[z,x]=p(x\mid z).
\]
The target is the posterior predictive:
\[
p(x\mid x_{1:t})
=
\sum_z P[z,x]\,p(z\mid x_{1:t}).
\]

Its logits can be computed without normalizing the intermediate posterior:
\[
\boxed{
\log\tilde p(x\mid x_{1:t})
=
\log\sum_z
P[z,x]\,\tilde p(z\mid x_{1:t}).
}
\]
The omitted posterior normalizer contributes only a constant across $x$. Centering this expression gives $L_c(s_t,x)$.

The logarithm after marginalization makes this kernel generally nonlinear in latent-inference coordinates. In particular, the log of a mixture need not lie in the span of the component log-distributions. The history kernel can therefore have higher rank than the centered static kernel, up to $|X|-1$.

When the posterior concentrates on a latent state and the component probabilities are positive, the predictive logits approach the corresponding static row. Their components outside the static logit span then vanish. The context-length curriculum consequently controls how strongly these uncertainty-dependent directions contribute to the Gram matrix; they are not necessarily low-variance or monotonically decreasing at every step.

The object fixed by next-token prediction is therefore the centered log-predictive, not the latent posterior. Histories with identical posterior predictives have identical minimal factorization codes even when their latent beliefs differ. This predicts a linear correspondence between readout-relevant activations and the log-predictive SVD code, while no analogous linear correspondence with the entire latent belief is guaranteed. A larger network may nevertheless retain additional belief information in directions invisible to the readout.

Both tasks therefore admit the same overall description:
\[
x_{1:t}
\;\xrightarrow{\Pi}\;
s_t
\;\xrightarrow{E}\;
E(s_t)
\;\xrightarrow{D^\top}\;
L_c(s_t,\cdot)^\top.
\]

Inference determines the conditional associated with each history. Kernel factorization determines its normative representation for linear readout. Different architectures can implement the history-to-state computation differently while being compared against the same kernel-factorization target.